\documentclass{article}
\usepackage{amsmath,amssymb,amsthm}
\usepackage{algorithm}
\usepackage{algpseudocode}
\usepackage{bm}
\usepackage{enumitem}
\newtheorem{theorem}{Theorem}
\newtheorem{lemma}{Lemma}
\newtheorem{assumption}{Assumption}
\newtheorem{definition}{Definition}
\newcommand{\E}{\mathbb{E}}
\newcommand{\R}{\mathbb{R}}

\begin{document}

\section*{FedProx with Local SGD}

\section*{Mathematical Formulation}
Consider $K$ participating clients.  
Each client $k\in\{1,\dots,K\}$ holds a local convex smooth function $f_k:\R^{d}\!\to\!\R$.  
The global objective is 
\[
F(\mathbf{x})\;:=\;\frac1K\sum_{k=1}^{K}f_k(\mathbf{x}) .
\]

Let $\mathbf{x}^{(t)}$ denote the \emph{global} model after the $t$‑th communication round.  
During round $t$ the server broadcasts $\mathbf{x}^{(t)}$ to all clients.  
Each client performs $E$ local SGD steps on the \emph{proximal} local sub‑problem
\[
\min_{\mathbf{x}\in\R^{d}}
\; \underbrace{f_k(\mathbf{x})}_{\text{local loss}}
\;+\;\frac{\mu}{2}\|\mathbf{x}-\mathbf{x}^{(t)}\|^{2},
\qquad \mu\ge 0,
\tag{1}
\]
where $\mu$ is the proximal coefficient.  

Denote by $\mathbf{x}^{(t)}_{k,0}:=\mathbf{x}^{(t)}$ the local copy at client $k$ before any update.
For $e=0,\dots,E-1$ the $e$‑th local SGD step on client $k$ is
\[
\mathbf{x}^{(t)}_{k,e+1}
\;=\;
\mathbf{x}^{(t)}_{k,e}
\;-\;
\eta\,
\bigl(
\nabla f_k(\mathbf{x}^{(t)}_{k,e};\xi^{(t)}_{k,e})
\;+\;
\mu\bigl(\mathbf{x}^{(t)}_{k,e}-\mathbf{x}^{(t)}\bigr)
\bigr),
\tag{2}
\]
where $\eta>0$ is the learning rate and $\xi^{(t)}_{k,e}$ denotes the random data sample (or mini‑batch) used to compute a stochastic gradient.  

After $E$ steps client $k$ sends $\mathbf{x}^{(t)}_{k,E}$ to the server, which aggregates by simple averaging:
\[
\mathbf{x}^{(t+1)}\;=\;\frac1K\sum_{k=1}^{K}\mathbf{x}^{(t)}_{k,E}.
\tag{3}
\]

\section*{Pseudocode}
\begin{algorithm}[H]
\caption{FedProx with Local SGD}
\begin{algorithmic}[1]
\State \textbf{Input:} learning rate $\eta$, proximal coefficient $\mu$, local steps $E$, total rounds $T$
\State Initialize $\mathbf{x}^{(0)}$ (e.g., $\mathbf{0}$)
\For{$t=0,\dots,T-1$}
    \State Server broadcasts $\mathbf{x}^{(t)}$ to all clients
    \For{each client $k$ \textbf{in parallel}}
        \State $\mathbf{x}^{(t)}_{k,0}\gets\mathbf{x}^{(t)}$
        \For{$e=0,\dots,E-1$}
            \State Sample $\xi^{(t)}_{k,e}$ and compute stochastic gradient $g^{(t)}_{k,e}\gets\nabla f_k(\mathbf{x}^{(t)}_{k,e};\xi^{(t)}_{k,e})$
            \State $\mathbf{x}^{(t)}_{k,e+1}\gets\mathbf{x}^{(t)}_{k,e}
                -\eta\bigl(g^{(t)}_{k,e}
                +\mu(\mathbf{x}^{(t)}_{k,e}-\mathbf{x}^{(t)})\bigr)$
        \EndFor
        \State Send $\mathbf{x}^{(t)}_{k,E}$ to server
    \EndFor
    \State Server updates $\displaystyle\mathbf{x}^{(t+1)}\gets\frac1K\sum_{k=1}^{K}\mathbf{x}^{(t)}_{k,E}$
\EndFor
\State \textbf{Output:} $\mathbf{x}^{(T)}$
\end{algorithmic}
\end{algorithm}

\section*{Dry Run Example}
Consider a single client ($K=1$) with scalar parameter $w\in\R$ and local loss
\[
f(w)=(w-3)^{2}.
\]
Set $w^{(0)}=0$, learning rate $\eta=0.1$, proximal coefficient $\mu=0$ (i.e., FedAvg), and $E=3$ local steps.

\begin{enumerate}[label=\textbf{Step \arabic*:}, leftmargin=*, nosep]
\item \textbf{Iteration $e=0$:}
    \[
    \nabla f(w^{(0)}) = 2(w^{(0)}-3)=2(0-3)=-6.
    \]
    Update:
    \[
    w^{(1)} = w^{(0)} - \eta\,\nabla f(w^{(0)}) = 0 - 0.1(-6)=0.6.
    \]
    Objective value $f(w^{(1)})=(0.6-3)^{2}=5.76$.

\item \textbf{Iteration $e=1$:}
    \[
    \nabla f(w^{(1)}) = 2(0.6-3) = -4.8.
    \]
    Update:
    \[
    w^{(2)} = 0.6 - 0.1(-4.8)=1.08.
    \]
    Objective $f(w^{(2)})=(1.08-3)^{2}=3.6864$.

\item \textbf{Iteration $e=2$:}
    \[
    \nabla f(w^{(2)}) = 2(1.08-3) = -3.84.
    \]
    Update:
    \[
    w^{(3)} = 1.08 - 0.1(-3.84)=1.464.
    \]
    Objective $f(w^{(3)})=(1.464-3)^{2}=2.3660.
    \]
\end{enumerate}
Thus after three local SGD steps the parameter has moved from $0$ to $1.464$, and the loss decreased from $9$ to $\approx2.37$.

\newpage
\section*{Assumptions}
\begin{assumption}[Convexity and Smoothness]\label{ass:conv-smooth}
For each client $k$ the local loss $f_k$ is convex and $L$‑smooth, i.e.
\[
f_k(\mathbf{y})\;\ge\;f_k(\mathbf{x})+\langle\nabla f_k(\mathbf{x}),\mathbf{y}-\mathbf{x}\rangle,
\qquad
\|\nabla f_k(\mathbf{x})-\nabla f_k(\mathbf{y})\|\le L\|\mathbf{x}-\mathbf{y}\|
\]
for all $\mathbf{x},\mathbf{y}\in\R^{d}$.
\end{assumption}

\begin{assumption}[Unbiased Stochastic Gradients]\label{ass:unbiased}
For any $k$, $e$, and $t$,
\[
\E_{\xi^{(t)}_{k,e}}\!\bigl[\nabla f_k(\mathbf{x}^{(t)}_{k,e};\xi^{(t)}_{k,e})\mid\mathbf{x}^{(t)}_{k,e}\bigr]
=
\nabla f_k(\mathbf{x}^{(t)}_{k,e}).
\]
\end{assumption}

\begin{assumption}[Bounded Variance]\label{ass:variance}
There exists $\sigma^{2}>0$ such that for all $k$, $e$, $t$,
\[
\E_{\xi^{(t)}_{k,e}}
\!\bigl[\|\nabla f_k(\mathbf{x}^{(t)}_{k,e};\xi^{(t)}_{k,e})-\nabla f_k(\mathbf{x}^{(t)}_{k,e})\|^{2}
\mid\mathbf{x}^{(t)}_{k,e}\bigr]
\le \sigma^{2}.
\]
\end{assumption}

\begin{assumption}[Bounded Dissimilarity]\label{ass:dissim}
There exists $B\ge0$ such that for all $\mathbf{x}$,
\[
\frac1K\sum_{k=1}^{K}\|\nabla f_k(\mathbf{x})-\nabla F(\mathbf{x})\|^{2}\le B^{2}.
\]
\end{assumption}

\begin{assumption}[Proximal Coefficient]\label{ass:mu}
The proximal parameter satisfies $0\le\mu\le L$.
\end{assumption}

\section*{Main Convergence Theorem}
\begin{theorem}[Convergence of FedProx with Local SGD]\label{thm:conv}
Let $\{\mathbf{x}^{(t)}\}_{t=0}^{T}$ be generated by the algorithm with fixed learning rate
\[
\eta \le \frac{1}{L+\mu}.
\]
Under Assumptions \ref{ass:conv-smooth}–\ref{ass:mu}, the averaged iterate
\[
\bar{\mathbf{x}}_{T}\;:=\;\frac1T\sum_{t=0}^{T-1}\mathbf{x}^{(t)}
\]
satisfies
\[
\E\!\bigl[F(\bar{\mathbf{x}}_{T})-F(\mathbf{x}^{\star})\bigr]
\;\le\;
\frac{\|\mathbf{x}^{(0)}-\mathbf{x}^{\star}\|^{2}}{2\eta T}
\;+\;
\frac{\eta L\sigma^{2}}{K}
\;+\;
\frac{\mu\eta L}{2}\,\|\mathbf{x}^{(0)}-\mathbf{x}^{\star}\|^{2}
\;+\;
\frac{\eta B^{2}}{2},
\tag{4}
\]
where $\mathbf{x}^{\star}\in\arg\min_{\mathbf{x}}F(\mathbf{x})$.
Consequently, choosing $\eta =\Theta\!\bigl(1/\sqrt{T}\bigr)$ yields the optimal rate
\[
\E\!\bigl[F(\bar{\mathbf{x}}_{T})-F(\mathbf{x}^{\star})\bigr]=\mathcal{O}\!\Bigl(\frac{1}{\sqrt{T}}\Bigr).
\]
\end{theorem}

\section*{Theoretical Properties}
\begin{itemize}
\item \textbf{Stability w.r.t. data heterogeneity.} The term $B^{2}$ captures client dissimilarity; a larger $B$ slows convergence, while $\mu>0$ mitigates this effect by pulling local iterates toward the global model.
\item \textbf{Effect of the proximal coefficient.} Increasing $\mu$ improves stability (the algorithm behaves more like a proximal point method) but introduces the additive term $\frac{\mu\eta L}{2}\|\mathbf{x}^{(0)}-\mathbf{x}^{\star}\|^{2}$ in the bound; thus $\mu$ should be chosen proportionally to $1/\sqrt{T}$ for optimal scaling.
\item \textbf{Comparison to FedAvg.} Setting $\mu=0$ recovers the classical FedAvg bound; the extra $\mu$‑dependent term explains empirically observed robustness of FedProx on heterogeneous data.
\item \textbf{Hyper‑parameter sensitivity.} The stepsize condition $\eta\le 1/(L+\mu)$ guarantees that the descent lemma holds for the proximal‑augmented local objective.
\end{itemize}

\section*{Discussion}
The theorem shows that, despite performing multiple local SGD steps, FedProx retains the $\mathcal{O}(1/\sqrt{T})$ convergence rate typical of stochastic convex optimization. The additional error terms are explicitly quantified in terms of the proximal coefficient $\mu$, the variance $\sigma^{2}$, and the client dissimilarity $B^{2}$. By tuning $\eta$ and $\mu$ according to the total number of communication rounds $T$, one can balance communication efficiency (large $E$) against statistical efficiency (small error constants).

\newpage
\section*{Preliminaries (Definitions and Lemmas)}
\begin{definition}[Proximal Local Objective]
For client $k$ at round $t$, define
\[
\Phi_{k}^{(t)}(\mathbf{x})\;:=\;f_k(\mathbf{x})+\frac{\mu}{2}\|\mathbf{x}-\mathbf{x}^{(t)}\|^{2}.
\]
\end{definition}
\begin{lemma}[Smoothness of $\Phi_{k}^{(t)}$]\label{lem:smooth-prox}
If $f_k$ is $L$‑smooth and $0\le\mu\le L$, then $\Phi_{k}^{(t)}$ is $(L+\mu)$‑smooth:
\[
\|\nabla\Phi_{k}^{(t)}(\mathbf{x})-\nabla\Phi_{k}^{(t)}(\mathbf{y})\|
\le (L+\mu)\|\mathbf{x}-\mathbf{y}\|,
\qquad\forall\mathbf{x},\mathbf{y}.
\]
\end{lemma}
\begin{proof}
\[
\nabla\Phi_{k}^{(t)}(\mathbf{x})
= \nabla f_k(\mathbf{x})+\mu(\mathbf{x}-\mathbf{x}^{(t)}).
\]
Subtracting the same expression at $\mathbf{y}$ and using $L$‑smoothness of $f_k$ gives
\[
\|\nabla\Phi_{k}^{(t)}(\mathbf{x})-\nabla\Phi_{k}^{(t)}(\mathbf{y})\|
\le \|\nabla f_k(\mathbf{x})-\nabla f_k(\mathbf{y})\|+\mu\|\mathbf{x}-\mathbf{y}\|
\le (L+\mu)\|\mathbf{x}-\mathbf{y}\|.
\tag*{$\square$}
\]
\end{proof}

\begin{lemma}[Descent Lemma for Stochastic Proximal SGD]\label{lem:descent}
Let $\mathbf{x}^{+}= \mathbf{x} -\eta\bigl(g+\mu(\mathbf{x}-\mathbf{x}^{(t)})\bigr)$ with $g$ an unbiased estimator of $\nabla f_k(\mathbf{x})$ and $\eta\le 1/(L+\mu)$. Then
\[
\E\bigl[\Phi_{k}^{(t)}(\mathbf{x}^{+})\mid\mathbf{x}\bigr]
\le
\Phi_{k}^{(t)}(\mathbf{x})
-\frac{\eta}{2}\|\nabla f_k(\mathbf{x})\|^{2}
+\frac{\eta^{2}(L+\mu)}{2}\sigma^{2}.
\tag{5}
\]
\end{lemma}
\begin{proof}
By Lemma~\ref{lem:smooth-prox} and the standard smoothness inequality,
\[
\Phi_{k}^{(t)}(\mathbf{x}^{+})
\le
\Phi_{k}^{(t)}(\mathbf{x})
+\langle\nabla\Phi_{k}^{(t)}(\mathbf{x}),\mathbf{x}^{+}-\mathbf{x}\rangle
+\frac{L+\mu}{2}\|\mathbf{x}^{+}-\mathbf{x}\|^{2}.
\tag*{[Step 1]}
\]
Insert the update definition:
\[
\mathbf{x}^{+}-\mathbf{x}
= -\eta\bigl(g+\mu(\mathbf{x}-\mathbf{x}^{(t)})\bigr).
\tag*{[Step 2]}
\]
Since $\nabla\Phi_{k}^{(t)}(\mathbf{x})=\nabla f_k(\mathbf{x})+\mu(\mathbf{x}-\mathbf{x}^{(t)})$, we obtain
\[
\langle\nabla\Phi_{k}^{(t)}(\mathbf{x}),\mathbf{x}^{+}-\mathbf{x}\rangle
=
-\eta\langle\nabla f_k(\mathbf{x})+\mu(\mathbf{x}-\mathbf{x}^{(t)}),\,g+\mu(\mathbf{x}-\mathbf{x}^{(t)})\rangle.
\tag*{[Step 3]}
\]
Expanding the inner product and using $\E[g\mid\mathbf{x}]=\nabla f_k(\mathbf{x})$ yields
\[
\E\bigl[\langle\nabla\Phi_{k}^{(t)}(\mathbf{x}),\mathbf{x}^{+}-\mathbf{x}\rangle\mid\mathbf{x}\bigr]
=
-\eta\|\nabla f_k(\mathbf{x})\|^{2}
-\eta\mu\|\mathbf{x}-\mathbf{x}^{(t)}\|^{2}.
\tag*{[Step 4]}
\]
The squared step norm satisfies
\[
\|\mathbf{x}^{+}-\mathbf{x}\|^{2}
= \eta^{2}\|g+\mu(\mathbf{x}-\mathbf{x}^{(t)})\|^{2}
\le 2\eta^{2}\bigl(\|g-\nabla f_k(\mathbf{x})\|^{2}
+\|\nabla f_k(\mathbf{x})+\mu(\mathbf{x}-\mathbf{x}^{(t)})\|^{2}\bigr).
\tag*{[Step 5]}
\]
Taking expectation and applying the variance bound (Assumption~\ref{ass:variance}) gives
\[
\E\bigl[\|\mathbf{x}^{+}-\mathbf{x}\|^{2}\mid\mathbf{x}\bigr]
\le 2\eta^{2}\bigl(\sigma^{2}+ \|\nabla f_k(\mathbf{x})\|^{2}+\mu^{2}\|\mathbf{x}-\mathbf{x}^{(t)}\|^{2}\bigr).
\tag*{[Step 6]}
\]
Plugging \([Step 4]\) and \([Step 6]\) into \([Step 1]\) and using $\eta\le 1/(L+\mu)$ (so that $1-\eta(L+\mu)\ge0$) yields
\[
\E\bigl[\Phi_{k}^{(t)}(\mathbf{x}^{+})\mid\mathbf{x}\bigr]
\le
\Phi_{k}^{(t)}(\mathbf{x})
-\eta\|\nabla f_k(\mathbf{x})\|^{2}
+\eta^{2}(L+\mu)\sigma^{2}
-\eta\mu\|\mathbf{x}-\mathbf{x}^{(t)}\|^{2}
+\eta^{2}(L+\mu)\mu^{2}\|\mathbf{x}-\mathbf{x}^{(t)}\|^{2}.
\tag*{[Step 7]}
\]
Since $\mu\le L$, the last two terms combine to a non‑positive quantity, which can be dropped, leaving
\[
\E\bigl[\Phi_{k}^{(t)}(\mathbf{x}^{+})\mid\mathbf{x}\bigr]
\le
\Phi_{k}^{(t)}(\mathbf{x})
-\frac{\eta}{2}\|\nabla f_k(\mathbf{x})\|^{2}
+\frac{\eta^{2}(L+\mu)}{2}\sigma^{2},
\tag*{[Step 8]}
\]
where we used $a-\frac{\eta}{2}a\le a-\eta a$ for $a=\|\nabla f_k(\mathbf{x})\|^{2}$. This proves the lemma.
\end{proof}

\section*{Main Proof (Step‑by‑Step Derivation)}
We now prove Theorem~\ref{thm:conv}. Let $\mathbf{x}^{\star}$ be a global minimizer of $F$.  
Define the global proximal objective
\[
\Psi^{(t)}(\mathbf{x})\;:=\;F(\mathbf{x})+\frac{\mu}{2}\|\mathbf{x}-\mathbf{x}^{(t)}\|^{2}.
\]

\begin{enumerate}[label=[Step \arabic*], leftmargin=*, nosep]
\item \textbf{One local step bound.}  
For any client $k$ and local step $e$, Lemma~\ref{lem:descent} applied to $\Phi_{k}^{(t)}$ gives
\[
\E\!\bigl[\Phi_{k}^{(t)}(\mathbf{x}^{(t)}_{k,e+1})\mid\mathbf{x}^{(t)}_{k,e}\bigr]
\le
\Phi_{k}^{(t)}(\mathbf{x}^{(t)}_{k,e})
-\frac{\eta}{2}\|\nabla f_k(\mathbf{x}^{(t)}_{k,e})\|^{2}
+\frac{\eta^{2}(L+\mu)}{2}\sigma^{2}.
\tag{6}
\]

\item \textbf{Summation over the $E$ local steps.}  
Taking total expectation and telescoping over $e=0,\dots,E-1$ yields
\[
\E\!\bigl[\Phi_{k}^{(t)}(\mathbf{x}^{(t)}_{k,E})\bigr]
\le
\Phi_{k}^{(t)}(\mathbf{x}^{(t)})
-\frac{\eta}{2}\sum_{e=0}^{E-1}\E\!\bigl[\|\nabla f_k(\mathbf{x}^{(t)}_{k,e})\|^{2}\bigr]
+\frac{E\eta^{2}(L+\mu)}{2}\sigma^{2}.
\tag{7}
\]

\item \textbf{Averaging over clients.}  
Define $\overline{\mathbf{x}}^{(t)}:=\frac1K\sum_{k=1}^{K}\mathbf{x}^{(t)}_{k,E}$ (the server model after round $t$).  
Using convexity of $f_k$ and Jensen’s inequality,
\[
F(\overline{\mathbf{x}}^{(t)})\le\frac1K\sum_{k=1}^{K}f_k(\mathbf{x}^{(t)}_{k,E}).
\tag{8}
\]

\item \textbf{Relating $\Phi_{k}^{(t)}$ to $F$.}  
By definition,
\[
\Phi_{k}^{(t)}(\mathbf{x}^{(t)}_{k,E})
= f_k(\mathbf{x}^{(t)}_{k,E})+\frac{\mu}{2}\|\mathbf{x}^{(t)}_{k,E}-\mathbf{x}^{(t)}\|^{2}.
\tag{9}
\]
Summing (9) over $k$ and using (8) gives
\[
\frac1K\sum_{k=1}^{K}\Phi_{k}^{(t)}(\mathbf{x}^{(t)}_{k,E})
\ge
F(\overline{\mathbf{x}}^{(t)})+\frac{\mu}{2K}\sum_{k=1}^{K}\|\mathbf{x}^{(t)}_{k,E}-\mathbf{x}^{(t)}\|^{2}.
\tag{10}
\]

\item \textbf{Bounding the proximal drift.}  
Since $\mathbf{x}^{(t+1)}=\overline{\mathbf{x}}^{(t)}$, the squared distance to the optimum satisfies
\[
\|\mathbf{x}^{(t+1)}-\mathbf{x}^{\star}\|^{2}
= \|\overline{\mathbf{x}}^{(t)}-\mathbf{x}^{\star}\|^{2}
\le
\frac1K\sum_{k=1}^{K}\|\mathbf{x}^{(t)}_{k,E}-\mathbf{x}^{\star}\|^{2}.
\tag{11}
\]

\item \textbf{Descent of the global distance.}  
For any $k$, expanding the norm yields
\[
\|\mathbf{x}^{(t)}_{k,E}-\mathbf{x}^{\star}\|^{2}
=
\|\mathbf{x}^{(t)}-\mathbf{x}^{\star}\|^{2}
+2\langle\mathbf{x}^{(t)}_{k,E}-\mathbf{x}^{(t)},\mathbf{x}^{(t)}-\mathbf{x}^{\star}\rangle
+\|\mathbf{x}^{(t)}_{k,E}-\mathbf{x}^{(t)}\|^{2}.
\tag{12}
\]
Taking expectations, summing over $k$, and using the unbiasedness of the stochastic gradients together with Lemma~\ref{lem:descent} (specifically the term $-\frac{\eta}{2}\|\nabla f_k\|^{2}$) leads to
\[
\E\!\bigl[\|\mathbf{x}^{(t+1)}-\mathbf{x}^{\star}\|^{2}\bigr]
\le
\|\mathbf{x}^{(t)}-\mathbf{x}^{\star}\|^{2}
-2\eta\E\!\bigl[F(\mathbf{x}^{(t)})-F(\mathbf{x}^{\star})\bigr]
+ \eta^{2}(L+\mu)\frac{\sigma^{2}}{K}
+ \mu\eta^{2}\|\mathbf{x}^{(t)}-\mathbf{x}^{\star}\|^{2}
+ \frac{2\eta B^{2}}{K}.
\tag{13}
\]
The term $\frac{2\eta B^{2}}{K}$ appears after applying Assumption~\ref{ass:dissim} to bound the variance of the client gradients around the global gradient.

\item \textbf{Rearranging (13).}  
Bring the objective gap to the left‑hand side:
\[
2\eta\E\!\bigl[F(\mathbf{x}^{(t)})-F(\mathbf{x}^{\star})\bigr]
\le
\|\mathbf{x}^{(t)}-\mathbf{x}^{\star}\|^{2}
-\E\!\bigl[\|\mathbf{x}^{(t+1)}-\mathbf{x}^{\star}\|^{2}\bigr]
+ \eta^{2}(L+\mu)\frac{\sigma^{2}}{K}
+ \mu\eta^{2}\|\mathbf{x}^{(t)}-\mathbf{x}^{\star}\|^{2}
+ \frac{2\eta B^{2}}{K}.
\tag{14}
\]

\item \textbf{Summation over communication rounds.}  
Sum \([Step 9]\) for $t=0,\dots,T-1$ and divide by $2\eta T$:
\[
\frac1T\sum_{t=0}^{T-1}\E\!\bigl[F(\mathbf{x}^{(t)})-F(\mathbf{x}^{\star})\bigr]
\le
\frac{\|\mathbf{x}^{(0)}-\mathbf{x}^{\star}\|^{2}}{2\eta T}
+\frac{\eta(L+\mu)\sigma^{2}}{2K}
+\frac{\mu\eta}{2T}\sum_{t=0}^{T-1}\E\!\bigl[\|\mathbf{x}^{(t)}-\mathbf{x}^{\star}\|^{2}\bigr]
+\frac{B^{2}}{K}.
\tag{15}
\]

\item \textbf{Bounding the averaged distance term.}  
Since the iterates are generated by a non‑expansive map (due to $\eta\le 1/(L+\mu)$), the distance sequence is non‑increasing, i.e.
\[
\E\!\bigl[\|\mathbf{x}^{(t)}-\mathbf{x}^{\star}\|^{2}\bigr]\le\|\mathbf{x}^{(0)}-\mathbf{x}^{\star}\|^{2},\quad\forall t.
\tag{16}
\]
Insert (16) into the third term of (15):
\[
\frac{\mu\eta}{2T}\sum_{t=0}^{T-1}\E\!\bigl[\|\mathbf{x}^{(t)}-\mathbf{x}^{\star}\|^{2}\bigr]
\le
\frac{\mu\eta}{2}\|\mathbf{x}^{(0)}-\mathbf{x}^{\star}\|^{2}.
\tag{17}
\]

\item \textbf{Final bound.}  
Combining (15)–(17) and noting that $\bar{\mathbf{x}}_{T}=\frac1T\sum_{t=0}^{T-1}\mathbf{x}^{(t)}$ together with convexity of $F$ gives
\[
\E\!\bigl[F(\bar{\mathbf{x}}_{T})-F(\mathbf{x}^{\star})\bigr]
\le
\frac{\|\mathbf{x}^{(0)}-\mathbf{x}^{\star}\|^{2}}{2\eta T}
+\frac{\eta L\sigma^{2}}{K}
+\frac{\mu\eta L}{2}\|\mathbf{x}^{(0)}-\mathbf{x}^{\star}\|^{2}
+\frac{\eta B^{2}}{2},
\]
which is exactly the bound (4) stated in Theorem~\ref{thm:conv}. ∎
\end{enumerate}

\section*{Rate Derivation (With Constants)}
Choosing $\eta = \frac{c}{\sqrt{T}}$ with $c\le \frac{1}{L+\mu}$ gives
\[
\frac{\|\mathbf{x}^{(0)}-\mathbf{x}^{\star}\|^{2}}{2\eta T}
= \mathcal{O}\!\Bigl(\frac{1}{\sqrt{T}}\Bigr),\qquad
\frac{\eta L\sigma^{2}}{K}= \mathcal{O}\!\Bigl(\frac{1}{\sqrt{T}}\Bigr),\qquad
\frac{\mu\eta L}{2}\|\mathbf{x}^{(0)}-\mathbf{x}^{\star}\|^{2}= \mathcal{O}\!\Bigl(\frac{1}{\sqrt{T}}\Bigr),
\]
and the term $\frac{\eta B^{2}}{2}$ is also $\mathcal{O}(1/\sqrt{T})$. Hence the overall convergence rate is $\mathcal{O}(1/\sqrt{T})$.

\section*{Special Cases}
\begin{itemize}
\item \textbf{Homogeneous data ($B=0$).} The bound reduces to the classical SGD rate with an extra $\mu$‑dependent term.
\item \textbf{No proximal regularization ($\mu=0$).} The theorem coincides with the convergence guarantee for FedAvg (local SGD) under the same assumptions.
\item \textbf{Deterministic gradients ($\sigma=0$).} The stochastic term disappears, yielding a deterministic rate $\mathcal{O}(1/T)$ after averaging.
\end{itemize}

\end{document}